% 第8章 8.1.4：ReLU 单元 rho(s)=max(s,0) 的有限组合。
% (a) f*(x)=|x-1/2|=rho(x-1/2)+rho(1/2-x)；
% (b) 结点 0,0.25,0.5,0.8,1、结点值 0.2,0.6,0.3,0.7,0.5 的线性样条，斜率 1.6,-1.2,4/3,-1；
%     f(x)=0.2+1.6 rho(x)-2.8 rho(x-0.25)+(38/15) rho(x-0.5)-(7/3) rho(x-0.8)，细线为逐项部分和 g_1,g_2,g_3。解析式绘制。
\begin{tikzpicture}[font=\figurefont\footnotesize]
\pgfplotsset{every axis/.style={axis lines=left,
  x axis line style={draw=FigureStroke,line width=.6pt},y axis line style={draw=FigureStroke,line width=.6pt},
  tick style={draw=FigureStroke},tick label style={font=\figurefont\scriptsize,text=FigureInk},
  title style={font=\figurefont\footnotesize},label style={font=\figurefont\scriptsize},
  legend style={draw=none,fill=none,font=\figurefont\scriptsize,cells={anchor=west}}}}
\begin{axis}[name=A,width=52mm,height=38mm,xmin=0,xmax=1,ymin=-0.02,ymax=0.6,xtick={0,0.5,1},ytick={0,0.25,0.5},
  title={(a) 两个单元表示 $f_*$},legend style={at={(0.5,-0.2)},anchor=north,legend columns=3}]
\addplot[draw=FigureGreen,line width=1.1pt,dashed,domain=0:1,samples=3] coordinates {(0,0) (0.5,0) (1,0.5)};
\addplot[draw=FigureYellow,line width=1.1pt,dash dot] coordinates {(0,0.5) (0.5,0) (1,0)};
\addplot[draw=FigureBlue,line width=1.6pt,opacity=.85] coordinates {(0,0.5) (0.5,0) (1,0.5)};
\legend{$\rho(x-\tfrac12)$,$\rho(\tfrac12-x)$,和}
\end{axis}
\begin{axis}[at={(A.east)},anchor=west,xshift=12mm,width=52mm,height=38mm,xmin=0,xmax=1,ymin=0,ymax=1.3,
  xtick={0,0.25,0.5,0.8,1},xticklabels={$0$,$t_1$,$t_2$,$t_3$,$1$},ytick={0,0.5,1},
  title={(b) 线性样条的 ReLU 分解},legend style={at={(0.5,-0.2)},anchor=north,legend columns=4}]
\addplot[draw=FigureMuted,line width=.6pt,domain=0:1,samples=2] {0.2+1.6*x};
\addplot[draw=FigureGreen,line width=.6pt,dashed] coordinates {(0,0.2) (0.25,0.6) (1,-0.3)};
\addplot[draw=FigureYellow,line width=.6pt,dash dot] coordinates {(0,0.2) (0.25,0.6) (0.5,0.3) (1,0.966667)};
\addplot[draw=FigureBlue,line width=1.1pt,mark=*,mark size=1.2pt,mark options={fill=FigureBlue}] coordinates {(0,0.2) (0.25,0.6) (0.5,0.3) (0.8,0.7) (1,0.5)};
\legend{$g_1$,$g_2$,$g_3$,$f$}
\end{axis}
\end{tikzpicture}
